\section{Real-Data Example: ADNI}
\label{sec:adni}

We complement the simulation study with a real-data example based on the
Alzheimer's Disease Neuroimaging Initiative (ADNI). The purpose of this section
is different from that of the simulations. The simulations are used for
controlled methodological evaluation, whereas the ADNI analysis is meant to
show that the borrowing problem also appears clearly in an observed dataset.

ADNI is not a randomized treatment study, so we do not treat this section as a
treatment-effect analysis. Instead, we use it as a cohort-difference risk
problem. The question is whether borrowing from an earlier ADNI cohort would
pull inference away from the later concurrent cohort, and whether the proposed
NPE formulation can better recover the concurrent risk level.

\subsection{ADNI setup}

We define the external cohort as ADNI1 and the concurrent cohort as the pooled
ADNIGO and ADNI2 samples. The analysis is restricted to subjects with mild
cognitive impairment (MCI) at baseline. The outcome is conversion to dementia
(or Alzheimer's disease) by approximately 24 months. For subject $i$, let
\[
Y_i =
\begin{cases}
1, & \text{if subject } i \text{ converts by approximately 24 months},\\
0, & \text{otherwise.}
\end{cases}
\]

For each subject, we use the baseline covariate vector
\[
X_i = (\text{age},\ \text{female},\ \text{education},\ \text{APOE4 count},\
\text{MMSE}_{\mathrm{bl}},\ \text{ADAS13}_{\mathrm{bl}},\
\text{CDRSB}_{\mathrm{bl}},\ \text{FAQ}_{\mathrm{bl}}).
\]

Table~\ref{tab:adni-cohort-summary} summarizes the two cohorts. The concurrent
cohort contains 376 subjects with an event rate of 0.1702, whereas the
external cohort contains 337 subjects with an event rate of 0.3769. This
difference is large enough to make naive borrowing questionable.

\begin{table}[t]
\centering
\caption{ADNI cohort summary for the real-data borrowing example.}
\label{tab:adni-cohort-summary}
\begin{tabular}{lcc}
\toprule
Cohort & Number of subjects & Conversion rate \\
\midrule
Concurrent (ADNIGO + ADNI2) & 376 & 0.1702 \\
External (ADNI1)            & 337 & 0.3769 \\
\bottomrule
\end{tabular}
\end{table}

\subsection{Baseline benchmark}

Before applying NPE, we first check whether the data show the borrowing issue
we are trying to study. We fit repeated logistic-regression baselines and
evaluate them on a held-out concurrent test set:
\begin{itemize}
    \item \textbf{ConcurrentOnly}: trained only on concurrent data;
    \item \textbf{ExternalOnly}: trained only on external data;
    \item \textbf{Pooled}: trained on the combined external and concurrent data.
\end{itemize}

The results are shown in Table~\ref{tab:adni-baselines}. All three approaches
have similarly high AUC and accuracy, so the main issue is not ranking
performance. The more noticeable difference is calibration. The
\texttt{ExternalOnly} and \texttt{Pooled} models both predict risks that are
too high relative to the concurrent cohort. The observed concurrent event rate
is 0.168, while the corresponding mean predicted risks are 0.248 and 0.218.
This suggests that direct borrowing mainly distorts the level of risk rather
than discrimination.

\begin{table}[t]
\centering
\caption{Baseline benchmark on a held-out concurrent test set.}
\label{tab:adni-baselines}
\begin{tabular}{lcccccc}
\toprule
Method & AUC & Accuracy & Brier & LogLoss & MeanPred & MeanObs \\
\midrule
ConcurrentOnly & 0.918 & 0.881 & 0.084 & 0.271 & 0.183 & 0.168 \\
ExternalOnly   & 0.923 & 0.887 & 0.088 & 0.303 & 0.248 & 0.168 \\
Pooled         & 0.923 & 0.883 & 0.083 & 0.276 & 0.218 & 0.168 \\
\bottomrule
\end{tabular}
\end{table}

\subsection{ADNI-specific NPE formulation}

Since ADNI does not provide the same randomized treatment setting as the
simulation study, we use a simpler formulation in which the target is the
concurrent risk level. Let $X_c$ and $X_e$ denote the observed covariates in
the concurrent and external cohorts. We simulate outcomes according to
\begin{align}
Y_c &\sim \mathrm{Bernoulli}\!\left(\sigma(X_c\beta + \theta)\right), \\
Y_e &\sim \mathrm{Bernoulli}\!\left(\sigma(X_e\beta + \theta + \delta)\right),
\end{align}
where $\sigma(z)=1/(1+e^{-z})$ is the logistic link.

In this model, $\theta$ is the concurrent risk-level parameter of interest,
$\delta$ represents the shift between the external and concurrent outcome
levels, and $\beta$ denotes the covariate effects. The role of $\delta$ is
important here because it allows the external cohort to differ rather than
forcing full exchangeability.

To generate training data for NPE, we first fit a logistic model to the
observed ADNI data and use the fitted coefficient vector as an anchor
$\beta_{\mathrm{anchor}}$. We then simulate outcomes on the real ADNI
covariates by sampling
\[
\theta \sim \mathrm{Uniform}(-3,0), \qquad
\delta \sim N(0,0.75^2),
\]
adding small perturbations around $\beta_{\mathrm{anchor}}$, and generating
$Y_c$ and $Y_e$ on the fixed covariate matrices $X_c$ and $X_e$. For each
simulated dataset, we compute summary statistics based on covariate means and
standard deviations, event-rate summaries, between-cohort differences, and
sample-size information. The NPE model is then trained to approximate the
posterior distribution of $\theta$.

\subsection{ADNI results}

Table~\ref{tab:adni-main-results} shows the main ADNI result. The posterior
mean of the concurrent risk parameter is $-1.5990$, with posterior standard
deviation $0.8585$ and a 95\% interval of $(-3.3009,\ 0.1082)$. The observed
concurrent event rate is 0.1702, which corresponds to an empirical concurrent
logit risk of
\[
\mathrm{logit}(0.1702) \approx -1.5841.
\]
The NPE posterior mean is therefore very close to the empirical concurrent
target.

\begin{table}[t]
\centering
\caption{Main ADNI real-data result for the NPE-based formulation.}
\label{tab:adni-main-results}
\begin{tabular}{lc}
\toprule
Quantity & Value \\
\midrule
Posterior mean of $\theta$     & $-1.5990$ \\
Posterior SD of $\theta$       & $0.8585$ \\
95\% interval for $\theta$     & $(-3.3009,\ 0.1082)$ \\
Observed concurrent event rate & $0.1702$ \\
Observed external event rate   & $0.3769$ \\
Empirical concurrent logit     & $-1.5841$ \\
\bottomrule
\end{tabular}
\end{table}

Table~\ref{tab:adni-comparators} compares this estimate with several simple
reference values. The empirical external-only logit and pooled logit are both
substantially above the concurrent target. Even the offset-based MLE estimates
remain noticeably higher than the concurrent empirical logit. Among the
estimators considered here, the NPE posterior mean is the closest to the
concurrent target.

\begin{table}[t]
\centering
\caption{Comparator estimates in the ADNI real-data example.}
\label{tab:adni-comparators}
\begin{tabular}{lc}
\toprule
Estimator & Value \\
\midrule
Empirical logit (concurrent)        & $-1.5841$ \\
Empirical logit (external)          & $-0.5029$ \\
Empirical logit (pooled)            & $-1.0054$ \\
Offset-MLE $\theta$ (concurrent)    & $-1.0749$ \\
Offset-MLE $\theta$ (pooled)        & $-0.7240$ \\
Offset-MLE $\delta$ (external shift)& $0.6022$ \\
NPE posterior mean $\theta$         & $-1.5990$ \\
\bottomrule
\end{tabular}
\end{table}

Overall, the ADNI analysis supports the same broad message as the simulation
study. The earlier and later cohorts are not directly comparable, and direct
borrowing tends to push the estimated risk upward. The ADNI-specific NPE
formulation is much closer to the concurrent target. We do not view this as a
replacement for the simulation study, but rather as a concrete example showing
that the borrowing problem is not only a theoretical concern.